% NORMA_COPPE_2026.tex -- versao composta da Norma da COPPE/UFRJ.
%
% A fonte de autoria e NORMA_COPPE_2026.md, na raiz do repositorio; este
% arquivo e a transcricao composta dela, mantida legivel como LaTeX e sem
% depender de pandoc. Editar os dois em conjunto.
%
% Usa a classe article de proposito: e um documento normativo SOBRE a classe
% coppe, nao uma tese, e aplicar a propria classe a ele confundiria.
%
% NAO sai do ufrj.dtx: a Norma tem ciclo de vida proprio, decidido pela CPGP,
% e nao e artefato da classe.

\documentclass[a4paper,11pt]{article}
\usepackage[T1]{fontenc}
\usepackage[utf8]{inputenc}
\usepackage{lmodern}
\usepackage[brazilian]{babel}
\usepackage[margin=2.5cm]{geometry}
\usepackage{titlesec}
\usepackage{enumitem}
\usepackage{microtype}
\usepackage{url}
\usepackage{parskip}
\usepackage{longtable}
\usepackage{booktabs}
\usepackage{array}
\usepackage[hidelinks]{hyperref}

\titleformat{\section}{\large\bfseries}{\thesection}{1ex}{}
\titleformat{\subsection}{\normalsize\bfseries}{\thesubsection}{1ex}{}

% Cada secao normativa declara de que tipo e a diferenca em relacao ao Manual.
\newcommand\tipo[1]{\normalfont\textit{--- #1}}

\title{\bfseries Norma para a Elaboração Gráfica de Teses e
       Dissertações da COPPE/UFRJ\\[2mm]
       \large Edição 2026 -- Proposta para apreciação da CPGP}
\author{}
\date{Setembro de 2026}

\begin{document}
\maketitle

\noindent\textit{Esta Norma, quando aprovada, substitui a} Norma para a
Elaboração Gráfica de Teses/Dissertações \textit{da COPPE/UFRJ aprovada pela
CPGP em 15 de julho de 2008 e suas revisões subsequentes (01/10/2009,
10/09/2010 e 26/11/2019).}

\hrulefill

\section*{Como ler esta Norma}

Esta Norma \textbf{não descreve como se formata uma tese}. Quem faz isso é o
\textbf{Manual para Elaboração e Normalização de Trabalhos Acadêmicos} da
UFRJ/SiBI, 9.\textsuperscript{a} edição revista (Rio de Janeiro, 2026) -- daqui
em diante \textit{Manual UFRJ 2026} --, que vale integralmente na COPPE.

O que esta Norma faz é registrar, ponto a ponto, \textbf{onde a COPPE
especializa o Manual}: onde ele deixa uma escolha em aberto e a COPPE escolhe,
onde ele pede um dado que só a COPPE tem, e onde a COPPE acrescenta um elemento
que ele não prevê. Fora do que está escrito aqui, vale o Manual --- formato do
papel, margens, espaçamento, fonte, paginação, equações, tabelas, indicativos de
seção, ortografia, siglas, ilustrações e estrutura textual no caso geral.

Cada seção abaixo diz de que tipo é a diferença:

\begin{description}[leftmargin=2.6cm,style=nextline,font=\bfseries]
  \item[Escolha] O Manual admite mais de uma forma; a COPPE fixa uma.
  \item[Dado próprio] O Manual pede a informação; o valor é da COPPE.
  \item[Acréscimo] Elemento que o Manual não prevê e a COPPE exige ou admite,
    onde o Manual não obriga e sem contrariar exemplo dele.
  \item[Reafirmação] Nada é acrescentado; repete-se um ponto do Manual por ser
    fonte frequente de erro.
\end{description}

\textbf{O que esta Norma pode decidir.} Ela \textbf{não contraria o Manual}:
onde o Manual obriga, vale o Manual. Ela decide \textbf{só onde o Manual não
obriga} --- onde ele deixa uma escolha em aberto ou pede um dado que só a COPPE
tem ---, e, ao decidir, \textbf{fica com o que os exemplos do Manual mostram}. As
duas folhas do modelo que o SiBI distribui --- a folha de rosto e a folha
adicional da Coleta CAPES, que a 3.1.2.1.2 manda baixar --- valem como
\textbf{errata} sobre os exemplos do Manual: onde as duas diferem, vale o
modelo. As normas da ABNT valem nas edições que o Manual cita, e por meio dele
(Seção~\ref{norma:citacoes}).

A edição revista de 2026 do Manual incorpora três decisões institucionais que
condicionam esta Norma: a \textbf{Resolução CEPG n. 246/2023}, que extingue a
via impressa e institui a entrega exclusivamente digital; a \textbf{Resolução
CEPG n. 302/2024}, cujo art.\ 57 fixa os idiomas de redação; e o \textbf{novo
modelo de coleta de informações acadêmicas da CAPES}, que dá origem à folha
adicional.

\subsection*{Mapa das diferenças}

\begin{center}
\begin{tabular}{@{}p{6.6cm}p{4.4cm}l@{}}
\toprule
\textbf{Nesta Norma} & \textbf{Especializa} & \textbf{Tipo}\\
\midrule
1. Logomarcas                          & Anexo A                & Escolha\\
2. Nome da instituição na capa         & 3.1.1                  & Escolha\\
3. Programas de Pós-graduação          & 3.1.1, 3.1.2.1.1       & Dado próprio\\
4. Identidade institucional em português & 3.1.1, 3.1.2.1.1     & Escolha\\
5. Orientação na folha de rosto        & 3.1.2.1.1(f)           & Escolha\\
6. Folha de resumo                     & 3.1.2.1.4, Anexos E e F & Reafirmação\\
7. Os resumos e o terceiro resumo     & 3.1.2, 3.1.2.1.4, 3.1.2.1.5 & Acréscimo\\
8. Referência no alto do resumo        & 3.1.2.1.4, Anexo E     & Escolha\\
9. Idioma principal                    & art.\ 57, CEPG 302/2024 & Escolha\\
10. Listas próprias                    & 3.1.2.2.4              & Acréscimo\\
11. Fonte das ilustrações              & 4.2                    & Escolha\\
12. Apêndices e Anexos                 & 3.1.4.2, 3.1.4.3       & Escolha\\
13. Numeração das folhas               & 2.7                    & Reafirmação\\
14. Sumário                            & 3.1.2.1.6              & Escolha\\
15. Citações e Referências             & 3.1.4.1                & Escolha\\
\bottomrule
\end{tabular}
\end{center}

\hrulefill

\section{Logomarcas \tipo{Escolha}}

O Anexo A do \textit{Manual UFRJ 2026} mostra, no alto da capa, dois logotipos,
e marca ambos como \textbf{opcionais}: o da UFRJ, à esquerda, e o do Programa,
à direita.

\textbf{A COPPE adota os dois}, e especializa o da direita: no lugar do logotipo
do Programa vai o do \textbf{Instituto (COPPE)}, o logotipo unificado, comum a
todos os Programas. A COPPE não usa o logotipo de cada Programa, até porque nem
todo Programa tem um. O da UFRJ, à esquerda, é fixado pela classe da UFRJ, e não
por esta Norma.

\begin{itemize}[nosep]
  \item Ambos figuram no alto da \textbf{capa}, na mesma linha, o da UFRJ
    alinhado à margem esquerda e o da COPPE à margem direita. \textbf{Não
    figuram na folha de rosto}: o Anexo A é o modelo da capa, e o Anexo B ---
    o da folha de rosto --- não traz imagem nenhuma, como não traz o modelo da
    folha de rosto que o SiBI distribui;
  \item Ambos com \textbf{2\,cm de altura}, preservada a proporção original;
  \item Nenhum outro logotipo é admitido --- nem de Programa, nem de
    laboratório, nem de agência de fomento. O vínculo com a agência é declarado
    na folha adicional da Coleta CAPES, não como imagem.
\end{itemize}

Os arquivos vetoriais em uso são distribuídos com a implementação de referência
(Seção~\ref{norma:implementacao}).

\section{Nome da instituição na capa \tipo{Escolha}}

O item 3.1.1 do Manual pede ``nome da instituição'', sem dizer em quantos
níveis. Na COPPE são \textbf{três linhas}, nesta ordem:

\begin{enumerate}[nosep]
  \item \textbf{Universidade Federal do Rio de Janeiro};
  \item \textbf{Instituto Alberto Luiz Coimbra de Pós-Graduação e Pesquisa de
    Engenharia};
  \item o nome do \textbf{Programa de Pós-graduação} em que o trabalho foi
    defendido.
\end{enumerate}

\textbf{Sem sigla na segunda linha.} O Anexo A do Manual, que é o modelo da
capa, escreve o nome da unidade por extenso e sem sigla (``ESCOLA DE BELAS
ARTES''). A sigla da COPPE não se perde: ela está no logotipo, no alto da mesma
folha. A frase da natureza da folha de rosto também traz o nome por extenso,
como o modelo do SiBI: ``\ldots ao Programa de Pós-Graduação em Engenharia de
Sistemas e Computação, Instituto Alberto Luiz Coimbra de Pós-Graduação e
Pesquisa de Engenharia, Universidade Federal do Rio de Janeiro, como requisito
parcial à obtenção do título de \ldots''.

\section{Programas de Pós-graduação da COPPE \tipo{Dado próprio}}

O Manual pede o nome do Programa; a lista é da COPPE. São treze, e o nome que
figura na capa, na folha de rosto e nas páginas de resumo é o da coluna
``Programa de Engenharia de\ldots''.

\begin{center}
\begin{tabular}{@{}llp{6.3cm}@{}}
\toprule
\textbf{Sigla} & \textbf{Programa} & \textbf{Nome em inglês}\\
\midrule
PEB  & Engenharia Biomédica                    & Biomedical Engineering\\
PEC  & Engenharia Civil                        & Civil Engineering\\
PEE  & Engenharia Elétrica                     & Electrical Engineering\\
PEM  & Engenharia Mecânica                     & Mechanical Engineering\\
PEMM & Engenharia Metalúrgica e de Materiais   & Metallurgical and Materials Engineering\\
PEN  & Engenharia Nuclear                      & Nuclear Engineering\\
PENO & Engenharia Oceânica                     & Ocean Engineering\\
PENT & Engenharia de Nanotecnologia            & Nanotechnology Engineering\\
PEP  & Engenharia de Produção                  & Production Engineering\\
PEQ  & Engenharia Química                      & Chemical Engineering\\
PESC & Engenharia de Sistemas e Computação     & Systems Engineering and Computer Science\\
PET  & Engenharia de Transportes               & Transportation Engineering\\
PPE  & Planejamento Energético                 & Energy Planning\\
\bottomrule
\end{tabular}
\end{center}

O nome em inglês é usado \textbf{apenas} na página de \textit{abstract}, onde o
Manual já prevê um resumo em língua estrangeira. Em toda a identidade
institucional o nome vai em português (Seção~\ref{norma:identidade}).

Alterações nesta lista --- criação, extinção ou renomeação de Programa --- são
comunicadas à equipe mantenedora da implementação de referência, e entram na
revisão seguinte do estilo da COPPE.

\section{Identidade institucional em português \tipo{Escolha}}
\label{norma:identidade}

O Manual não trata do caso de um trabalho redigido em idioma estrangeiro. Na
COPPE, os elementos de identidade institucional --- nome da Universidade, nome
do Instituto, nome do Programa, cidade (Rio de Janeiro), estado (RJ) e país
(Brasil) --- \textbf{figuram sempre em português}, qualquer que seja o idioma
principal do trabalho.

Acompanham o idioma principal apenas o \textbf{título} e o \textbf{subtítulo}.
A natureza e o objetivo ficam em português, como o resto das folhas de
identidade --- capa, folha de rosto e folha de aprovação ---, e a referência no
alto de cada resumo também (Seção~\ref{norma:refresumo}). A área de concentração
não entra na conta: ela é campo da folha adicional, que é um formulário da
Coleta CAPES em português, e não figura nas folhas de identidade.

\section{Orientação na folha de rosto \tipo{Escolha}}

O item 3.1.2.1.1(f) pede o nome do orientador e, quando houver, do
coorientador, sem fixar a posição. Na COPPE ambos figuram \textbf{alinhados a
partir da margem esquerda}, abaixo do bloco de natureza, o orientador antes do
coorientador.

\section{Folha de resumo \tipo{Reafirmação}}

Cada página de resumo segue o modelo dos Anexos E e F do Manual: o título da
folha (RESUMO, \textit{ABSTRACT}, RESUMEN), centralizado como todo título sem
indicativo numérico (2.6), a referência do trabalho
(Seção~\ref{norma:refresumo}), o texto do resumo e as palavras-chave.
\textbf{A COPPE não acrescenta nada a ela.}

Até esta edição, a COPPE acrescentava a frase de abertura (``Resumo da Tese
apresentada à COPPE/UFRJ\ldots''), o título com o nome do autor e o mês, e os
orientadores com o Programa. Nada disso está no modelo do Manual, e a
referência já traz o autor, o título, a natureza e a instituição. A
implementação de referência continua a oferecer esses elementos como opção,
desligada por padrão, para o Programa que os exija.

\section{Os resumos e o terceiro resumo \tipo{Acréscimo}}

O Manual prevê dois resumos, nesta ordem (3.1.2): o \textbf{resumo em língua
vernácula} (3.1.2.1.4), que é o português, e o \textbf{resumo em língua
estrangeira} (3.1.2.1.5), ``a versão do resumo em língua vernácula no idioma de
divulgação internacional'' --- o \textit{abstract} do Anexo F. \textbf{O resumo
em português vem sempre primeiro}, qualquer que seja o idioma do trabalho.

A COPPE fixa o inglês como a língua estrangeira e acrescenta um terceiro resumo
quando o trabalho não está escrito em nenhuma das duas línguas:

\begin{enumerate}[nosep]
  \item \textbf{Resumo em português} --- obrigatório, sempre o primeiro;
  \item \textbf{Resumo em inglês} --- obrigatório; num trabalho redigido em
    inglês, é o resumo no idioma do trabalho;
  \item \textbf{Resumo no idioma do trabalho} --- obrigatório só quando ele não
    é nem o português nem o inglês, e vem por último.
\end{enumerate}

A terceira posição existe por um caso concreto: num trabalho redigido em
espanhol, idioma admitido pelo art.\ 57 da Resolução CEPG n.\ 302/2024, os dois
resumos do Manual são o português e o inglês, e o resumo na língua em que o
trabalho foi escrito ficaria de fora. Ela vem depois dos dois do Manual, e não
entre eles.

Cada resumo ocupa uma página própria, na ordem acima, e encerra com as
palavras-chave \textbf{no seu próprio idioma}, separadas por ponto e vírgula,
como determina o 3.1.2.1.4.

\textbf{A página inteira fica no idioma do seu resumo} --- o título da folha, o
texto e as palavras-chave ---, com uma exceção: a referência no alto, que é a
mesma em todas as páginas e fica em português
(Seção~\ref{norma:refresumo}).

Isso vale só para as páginas de resumo. A capa, a folha de rosto e a folha de
aprovação são identidade institucional e continuam em português.

\section{Referência no alto do resumo \tipo{Escolha}}
\label{norma:refresumo}

O item 3.1.2.1.4 diz que ``\textit{sugere-se que o resumo venha antecedido por
uma referência conforme apresentado no ANEXO E}'', e os Anexos E e F mostram o
resumo e o \textit{abstract} começando por ela. \textbf{A COPPE adota a
sugestão como regra}: cada página de resumo abre pela referência do próprio
trabalho, logo acima do texto.

A referência segue o Anexo E e a NBR 6023:

\begin{quote}
SOBRENOME, Nome. \textbf{Título}: subtítulo. Rio de Janeiro, 2026. Tese (Doutorado em
Engenharia de Sistemas e Computação) --- Instituto Alberto Luiz Coimbra de
Pós-Graduação e Pesquisa de Engenharia, Universidade Federal do Rio de
Janeiro, Rio de Janeiro, 2026.
\end{quote}

É a \textbf{mesma} nas páginas de resumo e fica \textbf{sempre em português},
porque a referência descreve um documento e não se traduz --- é o que o Anexo F
faz ao repeti-la na página do \textit{abstract}. Só o título acompanha o
idioma do trabalho.

O \textbf{exame de qualificação} não a recebe, nem o \textbf{seminário de
mestrado} que alguns Programas exigem: nenhum dos dois é depositado na
biblioteca nem entra no acervo, e portanto nenhum dos dois é documento
referenciável.

A implementação de referência compõe a referência sozinha, com os dados que a
folha de rosto já exige, e oferece uma opção para retirá-la --- útil a quem
quer o resumo numa folha só, o que o Manual não exige.

\section{Idioma principal do trabalho \tipo{Escolha}}
\label{norma:idioma}

O trabalho pode ser redigido em \textbf{português}, \textbf{inglês} ou
\textbf{espanhol}, que são os idiomas admitidos pelo art.\ 57 da Resolução CEPG
n.\ 302/2024. Não há na UFRJ respaldo normativo para outros.

A escolha do idioma principal afeta o corpo textual, as legendas, os títulos das
seções pré-textuais e o título do trabalho na capa e na folha de rosto. Não
afeta a identidade institucional (Seção~\ref{norma:identidade}).

\section{Listas de Quadros, Programas e Algoritmos \tipo{Acréscimo}}

Em complemento às listas do item 3.1.2.2.4, a COPPE reconhece três listas
próprias, cada uma presente quando houver ao menos uma ocorrência no corpo do
trabalho. Quadro, programa e algoritmo são ilustrações (2.10), e por isso as
três vêm junto das demais listas de ilustrações, \textbf{antes} da Lista de
Tabelas, na ordem do item 3.1.2:

\begin{itemize}[nosep]
  \item \textbf{Lista de Quadros} --- quadro é a ilustração bordada em todos os
    lados, distinta da tabela, que é delimitada apenas no topo e na base;
  \item \textbf{Lista de Programas} --- listagens de código-fonte;
  \item \textbf{Lista de Algoritmos} --- pseudo-código.
\end{itemize}

\section{Fonte das ilustrações \tipo{Escolha}}

Para toda ilustração --- figura, tabela, quadro, programa, algoritmo --- a
legenda figura \textbf{acima} e a fonte é \textbf{obrigatória} e figura
\textbf{abaixo}, em corpo menor, espaçamento simples, antecedida da palavra
``Fonte:'' ou de seu equivalente no idioma principal. Legenda e fonte ficam
dentro das margens da ilustração (2.10) e alinhadas à esquerda, como na Figura~1
do Manual, o único exemplo que ele traz.

\section{Apêndices e Anexos \tipo{Escolha}}

Os capítulos de Apêndice e de Anexo são apresentados, tanto na abertura do
capítulo quanto no Sumário, no formato

\begin{quote}
  \textbf{APÊNDICE A --- Título do apêndice}\\
  \textbf{ANEXO A --- Título do anexo}
\end{quote}

com a palavra em caixa alta, a letra identificadora, travessão e o título,
respeitando o idioma principal.

\section{Numeração das folhas \tipo{Reafirmação}}

Nada se acrescenta aqui ao item 2.7 do Manual. A seção existe porque é o ponto
em que mais se observaram divergências de implementação.

A capa não é contada nem numerada. \textbf{A contagem começa na folha de rosto}
e prossegue por todas as folhas, que não são numeradas enquanto pré-textuais. A
folha adicional que contém a ficha catalográfica \textbf{não é contada nem
numerada}.

A numeração em algarismos arábicos é impressa \textbf{a partir da primeira folha
da parte textual, dando sequência à contagem já iniciada} --- e portanto a
Introdução \textbf{não} é a folha 1. Figura no canto superior direito, a 2\,cm
da borda superior, com o último algarismo a 2\,cm da borda direita.

Não se empregam algarismos romanos nas folhas pré-textuais.

\section{Sumário \tipo{Escolha}}

O item 3.1.2.1.6 manda alinhar os títulos à esquerda e \textbf{recomenda} que o
alinhamento seja ``pela margem do título do indicativo mais extenso'', mandando
usar como exemplo o sumário do próprio Manual. Por ser recomendação, a COPPE
fixa a forma:

\begin{itemize}[nosep]
  \item todo \textbf{indicativo} começa na margem esquerda, qualquer que seja o
    nível;
  \item todo \textbf{título} começa numa mesma coluna, cuja posição é dada pelo
    indicativo mais extenso do trabalho, separado dele por um espaço --- sem
    ponto, hífen ou travessão, como determina a 2.6;
  \item as linhas seguintes de um título longo alinham-se pela coluna do
    título;
  \item o pontilhado que liga o título ao número da folha corre em todos os
    níveis, e o espaçamento entre as linhas é uniforme.
\end{itemize}

A grafia de cada nível reproduz a do corpo do trabalho, como determina a 2.6, e
é a mesma que o Manual usa no seu próprio sumário:

\begin{center}
\begin{tabular}{@{}ll@{}}
\toprule
\textbf{Nível} & \textbf{Grafia}\\
\midrule
Primária    & CAIXA ALTA, negrito\\
Secundária  & CAIXA ALTA\\
Terciária   & negrito\\
Quaternária & negrito itálico\\
Quinária    & itálico\\
\bottomrule
\end{tabular}
\end{center}

Os elementos pós-textuais --- Referências, Apêndices, Anexos, Glossário,
Índice --- figuram no sumário com a grafia da seção primária.

\section{Citações e Referências \tipo{Escolha}}
\label{norma:citacoes}

Valem as normas da ABNT \textbf{nas edições que o Manual cita} na sua lista de
referências --- \textbf{NBR 6023:2025} (Referências) e \textbf{NBR 10520:2023}
(Citações) ---, e não ``as edições vigentes'', que podem mudar sem o Manual
mudar. Onde uma delas diverge do Manual, vale o Manual. Qualquer dos dois
sistemas de chamada que o Manual prevê --- autor-data ou numérico (4.1.1.1) ---
é admitido, a critério do autor, desde que usado de forma consistente em todo o
trabalho; o numérico não se usa em trabalho com notas de rodapé (4.1.1.1.1).

\hrulefill

\section{Implementação de referência}
\label{norma:implementacao}

A classe LaTeX \texttt{ufrj} com o estilo de unidade \texttt{ufrj-coppe},
versões 5.0 e posteriores (\textbf{CoppeTeX 5.x}), distribuída em
\url{https://github.com/COPPE-UFRJ/CoppeTeX}, é a implementação de referência
desta Norma; nas versões 4.0 e 4.1 foi a classe \texttt{coppe}. A classe
implementa o \textit{Manual UFRJ 2026}, e o estilo, o que esta Norma fixa como
dado próprio e como texto da COPPE. Produz, sem configuração adicional por parte
do autor além de carregar o estilo, um documento aderente ao mesmo tempo ao
\textit{Manual UFRJ 2026} e a esta Norma, e distribui os arquivos das duas
logomarcas.

Em caso de divergência entre o que está escrito aqui e o que a versão estável
corrente da classe produz, \textbf{prevalecem esta Norma e o Manual}, e a
divergência deve ser comunicada à equipe mantenedora para correção na revisão
seguinte.

\subsection{O que a implementação não faz}

\begin{itemize}[nosep]
  \item \textbf{A ficha catalográfica não é composta pela classe.} É obtida no
    Gerador de Fichas Catalográficas do SiBI
    (\url{http://fichacatalografica.sibi.ufrj.br/}) ou solicitada à biblioteca
    do Programa, e entra na folha adicional prevista no item 3.1.2.1.2,
    preenchida pelo Programa.
  \item \textbf{Os pacotes de francês e italiano distribuídos com a classe não
    habilitam esses idiomas para redação.} Servem de demonstração do mecanismo
    de extensão a novos idiomas; os idiomas admitidos são os da
    Seção~\ref{norma:idioma}.
\end{itemize}

\vspace{4mm}\hrulefill

\noindent\textit{Versão da proposta: setembro de 2026, revista para a
9.\textsuperscript{a} edição revista (2026) do Manual UFRJ/SiBI. A fonte de
autoria é} \texttt{NORMA\_COPPE\_2026.md}\textit{, na raiz do repositório.
\emph{Branch} \texttt{nlinguas} em}
\url{https://github.com/COPPE-UFRJ/CoppeTeX}\textit{.}

\end{document}
